\section{Настройка DNS-сервера BIND9}

\subsection{Настройка named.conf.options}

\begin{lstlisting}
nano /etc/bind/named.conf.options
\end{lstlisting}

\begin{lstlisting}
options {
    directory "/var/cache/bind";
    forwarders { 8.8.8.8; };
    dnssec-validation auto;
    auth-nxdomain no;
    listen-on { 127.0.0.1; 192.168.N.1; };
};
\end{lstlisting}

\texttt{forwarders}~--- DNS Google для внешних зон; \texttt{listen-on}~--- принимать запросы только с localhost и LAN-интерфейса.

\subsection{Настройка named.conf.local}

\begin{lstlisting}
nano /etc/bind/named.conf.local
\end{lstlisting}

\begin{lstlisting}
include "/etc/bind/rndc.key";
controls {
    inet 127.0.0.1 allow { localhost; } keys { rndc-key; };
};
zone "STUDENT.GROUP.local" IN {
    type master;
    file "/var/lib/bind/forward.db";
    allow-update { key rndc-key; };
};
zone "N.168.192.in-addr.arpa" IN {
    type master;
    file "/var/lib/bind/reverse.db";
    allow-update { key rndc-key; };
};
\end{lstlisting}

Файлы зон в \texttt{/var/lib/bind/}: здесь BIND имеет права на запись, что нужно для DDNS (журнальные файлы \texttt{.jnl}).

\subsection{Файл прямой зоны forward.db}

\begin{lstlisting}
nano /var/lib/bind/forward.db
\end{lstlisting}

\begin{lstlisting}
$TTL 86400
STUDENT.GROUP.local. IN SOA SERVERHOSTNAME.STUDENT.GROUP.local.
    admin.STUDENT.GROUP.local. (
        20110103 ; Serial
        604800   ; Refresh
        86400    ; Retry
        2419200  ; Expire
        604800 ) ; Negative Cache TTL
@       IN NS  SERVERHOSTNAME.STUDENT.GROUP.local.
@       IN A   192.168.N.1
localhost IN A  127.0.0.1
SERVERHOSTNAME IN A 192.168.N.1
\end{lstlisting}

\subsection{Файл обратной зоны reverse.db}

\begin{lstlisting}
nano /var/lib/bind/reverse.db
\end{lstlisting}

\begin{lstlisting}
$TTL 86400
N.168.192.in-addr.arpa. IN SOA STUDENT.GROUP.local.
    STUDENT.GROUP.local. (
        20110104 ; Serial
        10800    ; Refresh
        3600     ; Retry
        604800   ; Expire
        3600 )   ; Minimum
        IN NS  SERVERHOSTNAME.STUDENT.GROUP.local.
1       IN PTR STUDENT.GROUP.local.
1       IN PTR SERVERHOSTNAME.STUDENT.GROUP.local.
\end{lstlisting}

\subsection{Проверка и перезапуск BIND9}

\begin{lstlisting}
named-checkconf
named-checkzone STUDENT.GROUP.local /var/lib/bind/forward.db
named-checkzone N.168.192.in-addr.arpa /var/lib/bind/reverse.db
systemctl restart bind9
\end{lstlisting}

\subsection{Настройка netplan}

\begin{lstlisting}
nano /etc/netplan/00-installer-config.yaml
\end{lstlisting}

\begin{lstlisting}
network:
  ethernets:
    enp0s3:
      dhcp4: no
      addresses: [10.0.2.15/24]
      gateway4: 10.0.2.2
    enp0s8:
      dhcp4: no
      addresses: [192.168.N.1/24]
      nameservers:
        addresses: [192.168.N.1]
        search: [STUDENT.GROUP.local]
  version: 2
\end{lstlisting}

\begin{lstlisting}
netplan apply && reboot
\end{lstlisting}

\begin{lstlisting}
nslookup SERVERHOSTNAME
nslookup 192.168.N.1
\end{lstlisting}

\subsection{Настройка DDNS}

\textbf{Копирование ключа RNDC:}
\begin{lstlisting}
cp -Rr /etc/bind/rndc.key /etc/dhcp/ddns-keys/
\end{lstlisting}

\textbf{Обновление dhcpd.conf:}
\begin{lstlisting}
include "/etc/dhcp/ddns-keys/rndc.key";
ddns-updates on;
ddns-update-style standard;
ddns-domainname "STUDENT.GROUP.local";
zone STUDENT.GROUP.local. {
    primary 192.168.N.1; key rndc-key;
}
zone N.168.192.in-addr.arpa. {
    primary 192.168.N.1; key rndc-key;
}
authoritative;
subnet 192.168.N.0 netmask 255.255.255.0 {
    range 192.168.N.10 192.168.N.254;
    option domain-name-servers 192.168.N.1;
    option domain-name "STUDENT.GROUP.local";
    option routers 192.168.N.1;
    option broadcast-address 192.168.N.255;
    default-lease-time 604800;
    max-lease-time 604800;
}
\end{lstlisting}

\textbf{Перезапуск и проверка:}
\begin{lstlisting}
systemctl restart bind9
service isc-dhcp-server restart
nslookup DESKTOPNAME01
tail -40 /var/log/syslog
\end{lstlisting}
